\section*{September 7, 2026: Allowing price changes to depend on initial prices}
\textit{Codex draft; sensitivity check requested by Jacob.}

Allowing annual price changes to depend on each school area's initial price
level reduces the pooled dollar estimate with property controls from
\$56,983 to \$35,023, a 38.5\% reduction. The estimate remains positive,
but its 95\% interval runs from --\$8,439 to \$78,485. Initial price levels
therefore matter for the estimated contrast. This check does not establish
that the remaining increase is caused by closure or that the adjustment
removes a known amount of bias.

Initial price is each physical school site's mean sale price during
2008--2012, expressed in 2022 dollars and calculated from the same eligible
transactions used in the price analysis. The regression interacts this fixed
initial mean with every calendar year, omitting 2012. These interactions let
the dollar-price path vary linearly with initial price, with a separate common
slope in each year. They do not give every school site its own trend.
School-site fixed effects already absorb the initial-price main effect.

The sample, transaction weights, property controls, and clustering are
unchanged. The annual event studies use 10,921 sales at 29 treated and
48 control sites, with 2012 as reference. The pooled difference-in-differences
models use 9,912 sales because they exclude the closure year, 2013.
The comparison without property controls changes similarly:
\$62,880 becomes \$41,559, with an adjusted 95\% interval of
--\$5,955 to \$89,074.

\begin{center}
\begin{tabular}{lrr}
\hline
Pooled dollar estimate & Existing model & Initial price $\times$ year added \\
\hline
Site and year fixed effects & \$62,880 & \$41,559 \\
With property controls & \$56,983 & \$35,023 \\
\hline
\end{tabular}
\end{center}

The annual estimates also shrink while retaining an upward post-closure
pattern. With property controls, the 2016 estimate falls from \$77,474 to
\$41,905, the 2017 estimate from \$82,478 to \$51,596, and the 2018 estimate
from \$57,561 to \$31,979. Every adjusted annual treatment coefficient's
pointwise 95\% interval includes zero. The joint test of the four pre-closure
event coefficients has a p-value of 0.60 after adjustment, versus 0.23 before;
this is not evidence that the identifying assumption is established.

The initial-price covariate is estimated with uneven support. Median
pre-closure transaction counts are 40 per treated site and 65.5 per control
site. Four sites have fewer than ten pre-closure sales: Brown (2), Smyth (1),
Stewart (7), and Goldblatt (5). The initial means range from \$46,579 to
\$767,770 for treated sites and \$27,234 to \$721,859 for controls. Broad
range overlap does not guarantee that every treated site has a comparable
control. No sites are dropped in this check.

The adjustment is exploratory. Initial means partly reflect which properties
sold and which pre-closure years have more transactions. Because they are
constructed from pre-closure outcomes, measurement error and regression to
the mean are concerns. Reported standard errors condition on the constructed
initial-price variable. The relationship between initial prices and annual
changes is learned jointly from treated and control observations; it is a
modeling restriction, not an independently observed counterfactual. The
existing baseline remains available alongside this sensitivity specification.

{\small Reproduction: Make in the \texttt{home\_price\_distribution} audit
produces three new datasets (77 site rows, 82 coefficient rows, and eight model
rows), their standard reports, and the figure below. Existing source data and
sample definitions are unchanged; code is in the current working tree.
Unadjusted estimates and standard errors reproduce the baseline. Independent
OLS regressions with explicit site, year, and interaction indicators reproduce
the adjusted pooled and annual coefficients with property controls.}

\clearpage
\begin{center}\includegraphics[width=\textwidth]{../input/initial_price_events.png}\end{center}
Blue shows the existing dollar specification; red adds initial-price-by-year
interactions. Both panels use the same transactions. Intervals are pointwise,
clustered by school site, and conditional on the estimated initial-price
covariate. The vertical line marks 2013, and all event coefficients are relative
to 2012.
